
% ===
\section{Modelo Teórico}
% ===
% Definições >>> agentes >>> variáveis >>> ambiente produtivo >>> payoff >>> regimes ecológico-econômicos >>> proposições.

% (i) estrutura produtiva e fator ambiental;

% (ii) problema microeconômico estático do agente;

% (iii) definição estática dos regimes ecológico-econômicos;

% (iv) lei de movimento do fator ambiental;

% (v) problema intertemporal e equação de Bellman;

% (vi) estado estacionário comensal e condições de estabilidade;

% (vii) implicações agregadas para o modelo de Solow.


Lorem ipsum dolor sit amet, consectetur adipiscing elit. Sed do eiusmod tempor incididunt ut labore et dolore magna aliqua. Ut enim ad minim veniam, quis nostrud exercitation ullamco laboris nisi ut aliquip ex ea commodo consequat. Duis aute irure dolor in reprehenderit in voluptate velit esse cillum dolore eu fugiat nulla pariatur. Excepteur sint occaecat cupidatat non proident, sunt in culpa qui officia deserunt mollit anim id est laborum.

Preliminarmente, define-se o \textit{comensalismo econômico} como uma configuração ecológico-econômica na qual a atividade produtiva gera benefício líquido para o agente privado sem provocar deterioração líquida relevante das condições ambientais que participam da produção. Nessa configuração, a manutenção do fator ambiental integra a estratégia ótima do agente, em razão dos efeitos produtivos e dos incentivos institucionais incorporados ao seu cálculo econômico. Em sua forma dinâmica, o comensalismo corresponde a um estado estacionário no qual a viabilidade econômica da atividade produtiva e a manutenção do fator ambiental coexistem de maneira localmente estável.


